\begin{exercise}{21}
Let $\mathcal C$ be a category and $X$ a simplicial set. Recall that
$X^{\mathrm{op}}$ is the following simplicial set:
$X_n^{\mathrm{op}}=X_n$ and
$d_i^{\mathrm{op}}\colon X_n\to X_{n-1}$ is given by $d_{n-i}$,
and likewise $s_i^{\mathrm{op}}=s_{n-i}$ as a map $X_n\to X_{n+1}$.
Prove the following assertions:
\begin{enumerate}[label=(\arabic*),leftmargin=*,itemsep=0pt,parsep=0pt,
                  topsep=.5\baselineskip]
 \item $\mathrm N(\mathcal C^{\mathrm{op}})\cong
        \mathrm N(\mathcal C)^{\mathrm{op}}$
 \item $(\Delta^n)^{\mathrm{op}}\cong\Delta^n$
 \item $(\Lambda_i^n)^{\mathrm{op}}\cong\Lambda_{n-i}^n$
 \item $(\partial\Delta^n)^{\mathrm{op}}\cong\partial\Delta^n$
\end{enumerate}

\begin{proof}
Let $r_k\colon[k]\to[k]$ be the order-reversing set map
$r_k(t)=k-t$. For $\alpha\colon[a]\to[b]$ in $\Delta$, put
$R(\alpha)=r_b\alpha r_a$. This is order-preserving, and
$R(\beta\alpha)=R(\beta)R(\alpha)$ and $R^2=\mathrm{id}$.
It sends $d_i\colon[k-1]\to[k]$ to $d_{k-i}$ and
$s_i\colon[k+1]\to[k]$ to $s_{k-i}$. Thus the stated opposite
structure is $X^{\mathrm{op}}=X\circ R^{\mathrm{op}}$.

For (1), reverse each chain $F\colon[k]\to\mathcal C^{\mathrm{op}}$
to a chain $\Theta_k(F)\colon[k]\to\mathcal C$: its object at $i$
is $F(k-i)$, and its arrow $i\to j$ is the $\mathcal C$-arrow
underlying $F(k-j\to k-i)$. Reversal is a bijection, and for
$\beta\colon[\ell]\to[k]$ we have
\[
 \Theta_\ell(F\beta)=\Theta_k(F)R(\beta).
\]
This is precisely compatibility with the opposite simplicial
operators, proving (1) by Definition~1.1.43.

For (2), define in degree $k$
\[
 \phi_k\colon(\Delta^n)^{\mathrm{op}}_k\longrightarrow\Delta^n_k,
 \qquad \phi_k(\alpha)=r_n\alpha r_k.
\]
These are bijections, each inverse to itself. For
$\beta\colon[\ell]\to[k]$,
$\phi_\ell(\alpha R(\beta))=\phi_k(\alpha)\beta$, so they form
an isomorphism of simplicial sets.
Moreover, $\operatorname{im}\phi_k(\alpha)=r_n(\operatorname{im}\alpha)$.
Thus $\alpha$ omits some vertex different from $i$ exactly when
$\phi_k(\alpha)$ omits some vertex different from $n-i$, and
$\alpha$ is nonsurjective exactly when $\phi_k(\alpha)$ is.
By Definition~1.1.37, $\phi$ therefore restricts to the
isomorphisms in (3) and (4).
\end{proof}
\end{exercise}
